\begin{algorithm}\footnotesize

\caption{Kartographer}\label{algo:kino-map}

\KwIn{$\mathcal{R}$, the state of the robot; $discr$, discretizationSize}
\KwOut{$\mathcal{E}$, a set of edges approximating $\mathcal{F}$}

$\mathcal{E} \longleftarrow \emptyset$ \\
$q_{limits} \longleftarrow$  \textsc{getJointAnglesLimits}$(\mathcal{R})$  \\
$qs \longleftarrow$ \textsc{DiscretizeJointSpace($q_{limits}, discr$)} \\
\For{$q$ in $qs$}{
$ee_{pos} \longleftarrow$ \textsc{ForwardKinematics}$(q)$ \\
\If{\textsc{IsInWorkspaceLimits($ee_{pos}$)}}{
$e \longleftarrow Edge()$ \\
$e\cdot q_{init} \longleftarrow q$ \\
$stepSize \longleftarrow$ \textsc{SampleStepSize()} \\
$e\cdot q_{final} \longleftarrow q_{init}.Interpolate(stepSize)$ \\
$e\cdot qdot \longleftarrow$ \textsc{SamplePrimitiveVelocity()} \\
$\mathcal{E} \longleftarrow \mathcal{E} \cup e$
}
}


\Return{$\mathcal{E}$}

\end{algorithm}

% To generate
% the failure case edges for the robot, we first identify the
% new reachable volume of the robot in joint space, as shown
% in Figure 2. For the remaining useable joints, we identify
% principal motion joints that impart the most significant in-
% teraction point motion and make a dense grid of joint con-
% figurations between the lowest and highest ranges relevant
% to the environment. We apply dynamic motions to these
% principal motion joints, which provide the most movement
% and force in a specific direction and save their trajectories
% as edges.